\begin{RequAnswer}{4.1}
It means that the result of the algorithm is some value that can be assigned to a variable.
For instance, when a line of code like \verb+int a = min(4,5)+ executes,
the \verb+min+ method (algorithm) ``returns'' the minimum value so it can be assigned to
the variable \verb+a+.
\end{RequAnswer}
\begin{RequAnswer}{4.2}
Here it is. Note that the order of the parameters does not matter.
\begin{code}
double cuboidVolume(double w, double h, double d) {
    return w*h*d;
}
\end{code}
\end{RequAnswer}
\begin{RequAnswer}{4.3}
Here is one possibility: keep subtracting $n$ from $a$
until you get to a number less than $n$:
\begin{code}
int modN(int a, int n) {
    while(a>=n) {
        a=a-n;
    return a;
}
\end{code}
That is not very efficient if $a$ is large, however.
A more efficient solution would be as follows:
\begin{code}
int modN(int a, int n) {
    return a - (a/n)*n;
}
\end{code}
Make sure you understand why this solution works!
Also, it should be noted that both of these solutions require that $a\geq 0$.
\end{RequAnswer}
\begin{RequAnswer}{4.4}
As was the case there,
two reasonable solutions include \verb@(n+m/2)/m@ and \verb@(2n+m)/(2m)@.
\end{RequAnswer}
\begin{RequAnswer}{4.5}
Here is one possible solution:
\begin{code}
boolean congruentModN(int a, int b,int n) {
    if(a%n==b%n) {
        return true;
    } else {
        return false;
    }
}
\end{code}
Here is a very concise solution. Do you see why it works?
\begin{code}
boolean congruentModN(int a, int b,int n) {
    return((a-b)%n == 0);
}
\end{code}
\end{RequAnswer}
\begin{RequAnswer}{4.6}
The first one will {\em not} work if $a$ and $b$ have different signs.
The second one will always work since if the difference between
$a$ and $b$ is a multiple of $n$, it will always return $0$.
Thus, the second solution is better.
\end{RequAnswer}
\begin{RequAnswer}{4.7}
This one is pretty straightforward: \verb+if(x%2==1)+
\end{RequAnswer}
\begin{RequAnswer}{4.8}
The easy answer is as follows:
\begin{code}
if(x%2==0) {
   x++;
} else {
    x--;
}
\end{code}
One way to do it without a conditional statement is: \verb^x+=1-2*(x%2);^ Yikes!
\end{RequAnswer}
\begin{RequAnswer}{4.9}
It is probably a contingency.
A tautology is a proposition that is always true. Since there is a conditional statement involved, it is likely that different cases happen.
In other words, there would be no need for the \verb+if+ statement if the expression inside it is a tautology.
\end{RequAnswer}
\begin{RequAnswer}{4.10}
We need to use a nested conditional statement:
\begin{code}
if(!list.isEmpty()) {
    if(list.get(0)>= 100) {
        list.clear();
    }
}
\end{code}
\end{RequAnswer}
\begin{RequAnswer}{4.11}
\verb+if( x > 0 && y > 0 )+ (using DeMorgan's law). Note that
\verb+if( x >= 0 && y >= 0 )+ is {\bf not} correct since \verb+x>0+ is the negation of \verb+x<=0+.
\end{RequAnswer}
\begin{RequAnswer}{4.12}
$\verb+~+x = 218$   (11011010),
$x \verb+&+ y = 32$ (00100000),
$x \verb+|+ y = 117$ (01110101), and
$x \verb+^+ y = 85$ (01010101).
\end{RequAnswer}
\begin{RequAnswer}{4.13}
$n-1$. Notice that the loop starts at 1 and ends at $n-1$, so that's $n-1$ times
(not $n$ as you might have guessed).
\end{RequAnswer}
\begin{RequAnswer}{4.14}
Here is the solution I expected you to come up with:
\begin{code}
sumFirstN(int n) {
	int sum = 0;
	for( int i = 1; i <= n; i++) {
		sum+=i;
	}
}
\end{code}
O.K., I can't resist giving the better answer. The following is much more efficient, but it may
be unclear to you why it is correct. Do not worry--we will see the idea behind it later.
\begin{code}
sumFirstN(int n) {
	return n*(n + 1) / 2;
}
\end{code}
Do not worry if you do not understand this one---I
definitely did {\em not} expect you to come up with this one yet.
But I wanted to prime your mind for something we will learn very soon.
\end{RequAnswer}
\begin{RequAnswer}{4.15}
\begin{code}
int minimum(int a[], int n) {
    int min = a[0];
    for(int i=1;i<n;i++) {
        if(a[i]<min) {
            min=a[i];
        }
    }
    return min;
}
\end{code}
\end{RequAnswer}
\begin{RequAnswer}{4.16}
Does your solution use a \verb+boolean+ variable to keep track of whether or not you found a 0?
If so, try to rewrite your code so that it does not do so before you read the solution given below.
If you used a boolean variable, your solution is both more complicated and less efficient than it should be,
and it is best that you figure out the better way on your own so that you are less likely to write similar
code in the future.

Again, before reading the solution given below, does your solution contain an \verb+else+ clause?
If so, think about whether or not it is correct by walking through your code on a small array.
Since I can't read your code, I cannot be certain, but if you have an \verb+else+ clause,
your solution is likely to be incorrect.
\begin{code}
boolean containsZero(int a[], int n) {
    for(int i=0;i<n;i++) {
        if(a[i]==0) {
            return true;
        }
    }
    return false;
}
\end{code}
\end{RequAnswer}
\begin{RequAnswer}{4.17}
Here is the first version:
\begin{code}
boolean noZeroes(int a[], int n) {
    return !containsZero(a,n);
}
\end{code}
Here is the second version, which looks almost identical to the solution to the previous question:
\begin{code}
boolean noZeroes(int a[], int n) {
    for(int i=0;i<n;i++) {
        if(a[i]==0) {
            return false;
        }
    }
    return true;
}
\end{code}
\end{RequAnswer}
\begin{RequAnswer}{4.18}
Although these two conditional statements are
{\em logically equivalent}, they are actually not equivalent
in most programming languages due to short-circuiting.
The first one is correct because it makes sure
that the index is valid before indexing into the array.
The second one will probably crash if $i$ indexes outside of the array, so it is not correct.
For a more complete discussion, see Example~\ref{exa:ss}, since
it is addresses almost the same exact same question!
The only difference is that here we added a check to make sure the
index was not negative.
\end{RequAnswer}
\begin{RequAnswer}{4.19}
(a) The array has no entries that are 0.
(b) Here is the most obvious solution:
\begin{code}
boolean foo(int []A, int n)
    for (int i = 0; i < n; i++) {
        if(a[i] == 0) {
            return false;
        }
    }
    return true;
}
\end{code}
(c) A better name would be \verb+noZeroes+ or \verb+containsNoZeroes+
because it determine whether or not the array contains any zeroes, returning true if it contains no zeroes and
false otherwise.
\end{RequAnswer}
\begin{RequAnswer}{4.20}
(a) Translating it very literally, it is saying that it is not the case that there exists an element of the array that is zero.
In other words, the array contains no zeroes. In other words, this is saying the exact same thing as the expression
in Question~\ref{rq:fooBeforeBar}.
(b) The same algorithm from Question~\ref{rq:fooBeforeBar}, but replace \verb+foo+ with \verb+bar+.
(c) Same as answer to Question~\ref{rq:fooBeforeBar}.
\end{RequAnswer}
\begin{RequAnswer}{4.21}
\begin{code}
boolean allAre2MultipleOr3Multiple(int[] a, int n) {
   for(int i=0;i<n;i++) {
       if(a[i]%2!=0 && a[i]%3!=0) {
           return false;
       }
   }
   return true;
}
\end{code}
\end{RequAnswer}
\begin{RequAnswer}{4.22}
\begin{code}
int i=1;
while(i<n) {
    // do something
    i++;
}
\end{code}
\end{RequAnswer}
\begin{RequAnswer}{4.23}
Go through each element until you find a zero.
At the end, if we looked at the last element, there are
no zeros, so we return true. Otherwise, we return false.
\begin{code}
boolean noZeroes(int a[], int n) {
    int i=0;
    while(i<n && a[i]!=0) {
        i++;
    }
    if(i==n) {
        return true;
    } else {
        return false;
    }
}
\end{code}
\end{RequAnswer}
